\section{Solvation and hydration distributions}\label{sec:solvation}

Clusters in equilibrium with a vapor exchange solvent molecules. At a given
relative humidity, the hydrates of one core composition populate a Boltzmann
distribution over the solvent count. \CJKQC{} prints this distribution with
\flg{--hydration} for water and \flg{--solvation CODE} for any other solvent.

The distribution is built on the formation machinery of
Chapter~\ref{sec:formation}. A state's distribution value is its formation
free energy with the solvent's chemical potential set to its vapor pressure.
The distribution at a given relative humidity and the formation energy at a
given concentration are therefore two views of the same calculation.

\subsection{Flags and solvent pressure}\label{sec:solvation-flags}

\begin{center}
\begin{tabular}{@{}L{4.1cm}L{9.4cm}@{}}
\toprule
\textbf{Flag} & \textbf{Meaning} \\
\midrule
\flg{--hydration} & Print the raw distribution section for water, solvent
code \texttt{w}. \\
\flg{--hydration-corr} & Print the corrected distribution section for water;
requires corrections, attached or stored. \\
\flg{--solvation CODE} & Print the raw distribution section for the named
solvent code. \\
\flg{--solvation-corr CODE} & Print the corrected distribution section for
the named solvent code. \\
\flg{--rh VALUE} & Relative humidity as a fraction: \texttt{0.5} means 50
\%. Water only. \\
\flg{--psolvent VALUE} & Solvent partial pressure, using the \flg{--conc}
VALUE grammar. Valid for every solvent and overrides \flg{--rh}. \\
\bottomrule
\end{tabular}
\end{center}

All four section flags are print flags and each is a visible outcome on its
own. The plain spellings print the raw-level section; the \texttt{-corr}
spellings print the corrected-level section. The raw section prints before the
corrected one, whatever order the flags were given.

All solvent flags in one run must agree on a single solvent code. A second,
different code is a usage error, while repeated spellings of the same code are
redundant. Codes are non-empty, letter-only, case-significant strings.

\flg{--rh} is water physics: it must be a finite fraction greater than zero,
and values above 1 (supersaturation) are accepted. It is a usage error with a
solvent code other than \texttt{w}. \flg{--psolvent} accepts the same values
as \flg{--conc}: pascals, or a value suffixed with \texttt{ppt}, \texttt{ppb},
or \texttt{cmmc}.

The pressure resolution order is:

\begin{enumerate}
  \item \flg{--psolvent}, if given;
  \item otherwise \flg{--rh}, if given;
  \item otherwise the saturated vapor pressure, which is valid for water only.
\end{enumerate}

For a non-water solvent, neither \flg{--rh} nor the default applies: an
explicit \flg{--psolvent} is required.

For water, \flg{--rh} uses the Magnus correlation

\begin{equation}
p_{\mathrm{sat}}(T) = 610.94\,
  \exp\!\left(\frac{17.625\,\tau}{243.04+\tau}\right)\ \mathrm{Pa},
\qquad \tau = T - 273.15,
\end{equation}

which gives 610.94 Pa at 273.15 K and about $3.162\times10^3$ Pa at 298.15 K.
The resolved pressure $p_w = \texttt{rh}\cdot p_{\mathrm{sat}}(T)$ must be a
finite positive number before the run starts; a temperature at the domain edge
that makes it non-finite is a usage error.

The resolved pressure enters the run as a concentration setting supplied
automatically, as if \flg{--conc CODE p} had been given. The formation
sections' solvent term shifts accordingly, \flg{--cnt} scales it like any
concentration setting, and the distribution values carry it. Naming the
solvent code in an explicit \flg{--conc} is a usage error, because the
automatic setting already owns it; settings for other components ride along
untouched.

\subsection{States, cores, and distribution values}\label{sec:solvation-states}

The sections evaluate exactly the record set the formation sections evaluate:
the records that remain after filtering, every loaded record for a print-only
run, or the collapsed ensembles under \flg{--bavg}.

A \textbf{state} is one composition. Its \textbf{core} is that composition
without the solvent code. For example, with solvent \texttt{w}, the state
\texttt{1sa2w} has core \texttt{1sa}, and the states \texttt{1sa},
\texttt{1sa1w}, and \texttt{1sa2w} belong to the same core group. Records and
ensembles without components are excluded. States whose core is empty, such
as the pure solvent clusters \texttt{1w}, \texttt{2w}, and so on, are also
excluded: the solvent monomer itself prints no distribution row.

Core groups print in first-appearance order of their core among the printed
set. States within a group print in first-appearance order of their
composition.

A state's distribution value at the section's level is:

\begin{itemize}
  \item \textbf{uncollapsed}: the lowest free conformer, that is, the minimum
        over the state's members' defined level values. Members with an
        undefined value are ignored; a state with no defined member is
        undefined;
  \item \textbf{under \flg{--bavg}}: the ensemble's level value, because the
        collapse already provides one ensemble per composition. The
        collapse replaces the conformer minimum, including conformational
        entropy.
\end{itemize}

A state's formation value needs the monomer references of its full
stoichiometry. A missing or level-less reference code leaves the state
undefined at that level. A hydrated state whose solvent monomer was filtered
away prints a dash, while the solvent-free states of the same core can remain
defined.

\subsection{The distribution}\label{sec:distribution}

Within a core group, each defined state's weight is

\begin{equation}
w_{\mathrm{state}} = \exp\!\left(
  -\frac{v_{\mathrm{state}} - v_{\min}}{k_B T}\right),
\end{equation}

where $v_{\min}$ is the group's minimum defined value and $T$ is the run's
temperature. Its percentage is

\begin{equation}
100\,\frac{w_{\mathrm{state}}}{\sum_{\mathrm{defined}} w}.
\end{equation}

Undefined states contribute nothing and print \texttt{-}. A group with no
defined states prints all dashes. A group with exactly one defined state
prints 100.0.

The pressure physics rides inside the values. The implied solvent setting's term
$-n R T \ln(p_w/p^\circ)$ makes the weights equal the JKQC form
$(p_w/p^\circ)^n \exp(-(G-nG^\circ_{\mathrm{solvent}})/(k_B T))$. Every
non-solvent term common to the group cancels in the normalization.
\flg{--cnt}'s $(N-1)/N$ scaling varies with the state's total molecule count,
so under \flg{--cnt} the distribution simply follows the formation values.
The minimum subtraction keeps every exponent at or below zero, so no overflow
is reachable and a state far above the minimum underflows to zero without
affecting the rest.

\flg{--noname} never pools the core groups; the grouping is intrinsic to the
distribution. Sorting and \flg{--reverse} reorder the main table, never the
section, which keeps its own grouping order.

\subsection{Printed layout}\label{sec:solvation-layout}

There is no header row. For each core group, two rows are printed: the
states' composition tokens, then their percentages. Percentages print with
one decimal, \texttt{-} for an undefined state, and need not sum to exactly
100.0 because of rounding. A blank line follows every group's percentage row,
so groups read as blocks.

Columns align within a group: each column is as wide as the wider of its label
and percentage cell, cells are right-aligned, and there is a single space
between columns. For example:

\begin{verbatim}
 1sa 1sa1w 1sa2w
40.0  45.0  15.0

 1pinic 1pinic1w
 72.3    27.7
\end{verbatim}

The standard-output order is: main table, raw formation section, corrected
formation section, hydration/solvation sections (raw before corrected), then
the completion summary. The sections print over the same record set as the
formation sections, whether or not any formation section was requested.

\subsection{Solvation examples}\label{sec:solvation-example}

Print the water distribution at 50 \% relative humidity:

\begin{verbatim}
CJKQC hydrated.jsonl --hydration --rh 0.5 --filename --name
\end{verbatim}

Print both raw and corrected water distributions at 298.15 K:

\begin{verbatim}
CJKQC corrected.jsonl --hydration --hydration-corr --temp 298.15
\end{verbatim}

Print a sulfuric-acid solvation distribution at an explicit partial pressure
of $10^{10}$ molecules per cm$^3$:

\begin{verbatim}
CJKQC clusters.jsonl --solvation sa --psolvent 1e10cmmc
\end{verbatim}

Combine a concentration for another component with the water distribution:

\begin{verbatim}
CJKQC clusters.jsonl --hydration --rh 0.8 --conc am 10ppb
\end{verbatim}

Here the \flg{--conc am ...} setting shifts the formation values of the
ammonia component, while the automatic water setting carries the humidity; the
distribution normalizes within each core group.
